\section{Related Work}
\noindent \textbf{Time-efficient on-device training.}
A considerable amount of literature/frameworks have investigated how to enable and accelerate training on resource-constrained edge devices~\cite{huang2023elastictrainer, xu2022mandheling, 10.1145/3295500.3356156, mercury2021sensys, 10.1145/3384419.3430716, zhu2023pockengine, tflite, tfjs, MNN}.
For instance, ElasticTrainer~\cite{huang2023elastictrainer} reduces training time by non-deterministically training a portion of the network.
Mandheling~\cite{xu2022mandheling} directly downsamples the entire network and maps it to the hardware accelerator.
There are also many open-source frameworks tailored for on-device training, such as TFLite~\cite{tflite} from Google and MNN~\cite{MNN} from Alibaba.
\sys is compatible with the above frameworks and hardware accelerator oriented works.
The difference between \sys and other works, such as ElasticTrainer and PruneTrain, which also employ dynamic neural network architectures, lies in \sys not compromising on parameter efficiency.

\noindent \textbf{Neural network compression/approximation.}
Neural network compression/approximation refers to approximating the original network with simpler and lighter networks, which can come from network structure approximation (e.g., structural pruning~\cite{torchpruninggithub, krishnan2019structural, AdaSpring, AdaEnlight} or sparse weights~\cite{tangandyang2023torchsparse}) or data precision approximation (e.g., INT8/INT4 quantization~\cite{frantar2023gptq, Zafrir_2019}).
% Although these works can reduce the over-parameterization on a certain downstream task and statically generate a more compact network, they compromise the generality on other tasks.
Different from the aforementioned works, \sys's NN approximation is dynamic for each data point, and is accurate enough for lossless training.

\noindent \textbf{Foundation models as a mobile system service.}
Compared to traditional on-device AI paradigm that individually deploys a compact model for each task~\cite{yin2024llmservicemobiledevices, cai2018proxylessnas, cai2018proxylessnas, Fang_2018, lu2024smalllanguagemodelssurvey, xu2024surveyresourceefficientllmmultimodal, yin2024elmselasticizedlargelanguage}, recent advances in LLMs~\cite{zhou2023large, muennighoff2023crosslingual, BELLE, koala_blogpost_2023, brown2020language} and multimodal~\cite{tang2023anytoany, su2023pandagpt, radford2021learning, li2021align, li2022blip, li2023blip2} suggest that a highly parameterized 
OS-level foundational model~\cite{yuan2023rethinking, yang2023edgefm} with stronger generalization and emergent capabilities is more suitable for the evolving workloads in on-device AI.
\sys sheds light on building such a one-for-all foundation ML system on devices to continuously and personally serve ubiquitous AI tasks.

\noindent \textbf{Routing diverse data to multiple networks.}
\yws{
Several prior arts~\cite{tabi, li2021appealnetefficienthighlyaccurateedgecloud, kag2023efficient} also leverage the idea of routing diverse data to multiple networks that has diverse compute cost.
Appealnet~\cite{li2021appealnetefficienthighlyaccurateedgecloud} proposes a two-head little network for routing data to two cascaded big-little networks at inference time.
The difficulty in Appealnet is defined as a loose measurement on softmax probability.
Selective Query~\cite{kag2023efficient} trains a routing model for assigning the to-be-inferenced sample to either the cloud or the device.
These work mainly focus on \textit{inference} optimizations, which are much more simple and deterministic compared to training.
\sys's data routing identifies several unique aspects in training, including feature-level difficulty and training sample importance to convergence.
}

\section{Discussion}
\label{subsec: discussion}

\textbf{Data diversity.}
\yws{
\sys is the first work to identify and orchestrate data difficulty and importance to dynamically approximate the frozen layers of PET.
Although there are some early efforts~\cite{teerapittayanon2017branchynetfastinferenceearly, baldock2021deeplearninglensexample} that explore the diversity, their definition cannot be directly applied to PET mainly due to the following reasons.
Firstly, they mainly focus on reducing the inference-time overhead by layer-level early exit for diverse samples.
NN inference has a higher tolerance for extracted features.
For instance, when a feature is tortured, one classification network can still output the correct answer if the highest logit is in the label class, regardless the other classes.
While in training, this may lead to a severely affected convergence.
Instead, \sys sets stringent restriction for training by a feature-level matching in its difficulty definition.
Secondly, early exiting is essentially a pruning method with sequential layer-level heuristic, which is much more coarse-grained than \sys's IDC pruning that prunes neurons in independent dimensions.
}
% \mwx{(1) SOTA data diversity definition methods. (2) their drawbacks (3) Ours's difference.}

\noindent \textbf{Handling diversity distribution drift.}
\yws{
One practical issue that \sys faces in real-world deployment is the data diversity distribution drift.
As \sys focuses on on-device fine-tuning, the distribution of on-device data may differ from on-cloud pre-training data.
To this end, \sys has the following conclusions and alternatives.
Firstly, as elaborated in $\S$\ref{subsec: sensitivity} and $\S$\ref{subsubsec:policy model}, the pattern of data difficulty is much more simple and deterministic than the distribution on concrete tasks, thus showing strong transferablity.
For instance, an NN trained on ImageNet still needs fine-tuning on caltech-101 dataset, while the difficulty pattern that captured by \sys's policy model on ImageNet has been nearly equivalent to caltech-101.
Besides, \sys has several alternatives when the fine-tuning data is completely different from pre-training.
One can crowd-sourcing a portion of user-permitted trainable data to the cloud, federated learning the policy model, or directly on-device continuously training the policy model.
}

\noindent \textbf{Other approximation techniques.}
\yws{
Here we discuss the design philosophy behind \sys on choosing the approximation technique.
Sparsity and quantization is not a good fit for \sys, as they can only provide a limited wall-clock speedup on mobile SoC processors.
For instance, to the best of our knowledge, the Raspberry Pi has no dedicated compute units for quantized data formats like int4.
Thereby, a highly quantized network only achieves an insignificant speedup on these devices since most of its operators are compute-bound.
Sparsity is also not well-supported on mainstream mobile SoCs.
Early exiting is essentially a coarse-grained pruning, thus also sub-optimal compared to pruning.
\sys employs structural pruning, which is the most suitable technique direction for PET acceleration on mobile devices.
}

\noindent \textbf{Applicability to language models.}
\yws{
Currently, \sys is mainly designed for on-device tuning CV models.
The LLMs are typically much heavier than CV models (e.g., llama-7B), and their fine-tuning on extremely resource-constrained mobile devices is still not a practical solution for demands like personalization.
Nevertheless, the main techniques of \sys are generally applicable to LLMs, with minor adjustments like supporting LoRA structures and profiling token-level data diversity.
We leave this as our future work.
}